% 第5章图：四变量二值链的递推表。数据：phi1=(0,3),phi2=(2,0),phi3=(3,0),phi4=(1,2)，相邻取值不同成本2。
% D1=(0,3), D2=(2,2), D3=(5,2), D4=(5,4)；回溯表 p2=(0,0), p3=(0,1), p4=(1,1)；最优 (0,1,1,1)，成本4。
\begin{tikzpicture}[x=1mm,y=1mm,>=Stealth,
  cell/.style={rectangle,rounded corners=1.8pt,draw=FigureStroke,fill=white,line width=.7pt,minimum width=12mm,minimum height=7mm,inner sep=1pt},
  best/.style={cell,draw=FigureGreen,fill=FigureGreenFill,line width=.9pt},
  note/.style={font=\FigureNoteFont\fontsize{7.5}{9}\selectfont,text=FigureMuted}]
\def\cx{28}
% cells: names c<k><s>
\foreach \k/\a/\b/\sa/\sb in {1/0/3/best/cell,2/2/2/cell/best,3/5/2/cell/best,4/5/4/cell/best}{
  \node[\sa] (c\k0) at ({(\k-1)*\cx},0) {$\a$};
  \node[\sb] (c\k1) at ({(\k-1)*\cx},-13) {$\b$};
}
\foreach \k/\u in {1/{(0,3)},2/{(2,0)},3/{(3,0)},4/{(1,2)}}{
  \node[font=\figurefont\footnotesize] at ({(\k-1)*\cx},12) {$D_{\k}$};
  \node[note] at ({(\k-1)*\cx},7.2) {$\phi_{\k}=\u$};
}
\node[note,anchor=east] at (-8,0) {$x_k=0$};
\node[note,anchor=east] at (-8,-13) {$x_k=1$};
% all candidates (gray)
\foreach \k [evaluate=\k as \j using int(\k+1)] in {1,2,3}{
  \foreach \s in {0,1}{\foreach \t in {0,1}{
    \draw[FigureGrid,line width=.6pt] (c\k\s.east)--(c\j\t.west);
  }}
}
% selected predecessors (dark), backtrack chain (green)
\foreach \from/\to in {10/20,20/30,31/40}{\draw[->,FigureStroke,line width=.8pt] (c\from.east)--(c\to.west);}
\foreach \from/\to in {10/21,21/31,31/41}{\draw[->,FigureGreen,line width=1.1pt] (c\from.east)--(c\to.west);}
\node[note,align=center] at ({1.5*\cx},-23) {灰线为全部候选，箭头为选中的前驱；相邻取值不同时成对成本为$2$};
\end{tikzpicture}
